% !TeX root = ../../main.tex
\chapter{面向计算范围冗余的结构感知选择性计算与协同加速方法}\label{chap:paretoes}
\section{研究动机}
\subsection{全量候选扫描的计算冗余}
第\ref{chap:acceles}章通过低比特粗筛和高精度重算降低每个候选的平均处理成本，但第一阶段依然访问全部矩阵行。数据库规模扩大后，即使每个记录已经较为紧凑，完整扫描仍占据主要时间。对于具有候选相关性结构的检索数据，查询结果往往集中于少量区域，遍历所有区域并非达到给定召回目标的必要条件。

全量扫描的另一个限制是缺少面向质量目标的灵活调节。当应用允许降低部分召回以改善延迟时，固定扫描范围仍付出近似不变的成本。若能将候选空间划分为可独立探测的簇，并依据查询选择相关簇，就可以将质量需求直接映射为实际访问量。

\subsection{选择性计算的性能优化机会}
ParetoES 在低比特稀疏执行基础上引入簇级选择性计算\cite{paretoes2026}。离线阶段构建与稀疏向量方向相符的索引，在线阶段先比较查询与簇代表，再对选中簇内的候选执行相似度计算。探测簇数 $n_{\mathrm{probe}}$ 控制计算预算，使系统能够在多个召回目标下工作。

本章以召回率--吞吐率曲线刻画质量与成本的权衡。在测量的配置集合中，若一个配置在召回率和吞吐率上均不低于另一个配置，且至少一项严格更优，则前者支配后者。非支配配置构成经验 Pareto 前沿，用于选择满足给定召回目标的高吞吐工作点。

\subsection{选择性计算引入的不规则执行问题}
选择性执行改变了工作负载的组织方式。首先，通用聚类对稀疏向量的方向结构未必适配，簇代表与成员之间的偏差可能增加为达到目标召回必须扫描的范围。其次，簇大小和访问热度不同，使原本较规则的全量分区出现查询相关的负载偏斜。最后，在线簇筛选与候选排序引入额外比较逻辑，集中式排序电路可能占用过多资源并拉长关键路径。
\thesisfig[1.00]{paretoes/problem_describe.pdf}{稀疏聚类与热点访问引入的选择性计算挑战}{fig:pareto-problem}{%
图\ref{fig:pareto-problem}概括上述现象。为此，本章联合设计方向敏感聚类、硬件簇探测、局部排序与多核负载分配，将候选范围缩减转化为关键路径上的执行时间下降。
}


\section{结构感知选择性计算系统设计概述}
ParetoES 包含离线索引构建、数据准备和在线检索。离线阶段执行 Spherical K-means++ Refine，生成候选簇及稀疏簇代表；随后量化、剪枝并按簇重排矩阵，结合代表查询构建热点统计，生成簇到计算核的映射。在线阶段由 ACPE 并行完成簇代表比较、目标簇定位、簇内相似度计算和候选输出。
\thesisfig[1.00]{paretoes/SystemOverview.pdf}{ParetoES 离线构建与在线选择性检索流程}{fig:pareto-overview}{%
图\ref{fig:pareto-overview}中，索引和执行共享簇粒度，使候选选择结果能够直接转化为 HBM 访问区间。簇代表与所属候选按计算核共同放置，减少跨核寻址。参数通过主机配置传递，运行时调整探测预算无需重新综合 FPGA。该设计适用于索引可在多次查询中复用的场景，频繁数据更新则需额外考虑重聚类和数据迁移。
}


表\ref{tab:inheritance}明确本章与第三章的关系。INT6 和 Ultra-CSR 继续承担压缩与流式执行基础，ReSparse 按聚类检索需求进行调整。新增机制集中于候选簇选择、选择性执行和热点均衡，因此二者是从降低单位成本到缩减计算范围的技术递进。
\begin{table}[htbp]\centering
\caption{ParetoES 对 AccelES 技术基础的继承与拓展}\label{tab:inheritance}
\begin{tabularx}{\textwidth}{p{3cm}XX}\toprule
技术要素&AccelES&本章处理方式\\\midrule
低比特数据&INT6 全量筛选&沿用低位宽，优化归一化与量化流程\\
稀疏传输&Ultra-CSR 连续扫描&沿用编码，按簇组织与选择性读取\\
非零剪枝&全量候选上的 ReSparse&与簇探测联合校准，调整统计阈值\\
候选范围&全矩阵低比特扫描&新增方向敏感聚类与簇探测\\
核心架构&Q-core 与 F-core 分层计算&新增 ACPE 及分布式微型排序\\
核间分配&行分区与随机重排&新增基于热点统计的\hbalance{}\\\bottomrule
\end{tabularx}
\end{table}

\section{面向计算范围优化的候选选择方法}
\subsection{稀疏数据的聚类特征分析}
对于相似度检索，向量长度和方向可能承载不同含义。直接以欧氏距离聚类时，幅值差异可能主导划分，使方向相近但长度不同的向量分离。稀疏向量均值又可能在多个成员的非零支撑集上产生较密集的数值，从而增加簇代表计算成本，并与原始稀疏成员分布不一致。
\thesisfig[0.85]{paretoes/kmeans_vs_spherical.pdf}{欧氏距离聚类与方向敏感聚类的示意对比}{fig:pareto-cluster}{%
本章对向量进行范数归一化，按内积衡量方向相似性，并选取接近平均方向的实际稀疏成员作为簇代表。实际成员保留原始稀疏结构，可避免直接使用稠密均值带来的访问负担。图\ref{fig:pareto-cluster}是低维投影示意，其作用是解释几何差异；聚类有效性由第\ref{chap:analysis}章的召回曲线及本章实验验证。
}


\subsection{Spherical K-means++ Refine 聚类方法}
令归一化向量为 $\bar x_i=x_i/\norm{x_i}_2$，当前代表集合为 $\mathcal{M}$。初始化使用方向距离进行概率采样：
\begin{equation}
 p(\bar x_i)=\frac{1-\max_{\mu\in\mathcal M}\bar x_i^{\mathsf T}\mu}
 {\sum_j(1-\max_{\mu\in\mathcal M}\bar x_j^{\mathsf T}\mu)}.
 \label{eq:spherical-init}
\end{equation}
该过程提高初始代表的方向多样性。原实现从大小为 $\min(0.01m,10000)$ 的候选子集中采样，以控制大规模初始化成本。零向量和分母为零的退化情况需要在预处理阶段单独处理，避免将无定义余弦值交给后续计算。

迭代阶段先按最大内积将成员分配到簇，再以与归一化平均方向最相近的成员更新代表：
\begin{equation}
 \mu_c=\argmax_{x\in\mathcal C_c} x^{\mathsf T}
 \frac{\sum_{u\in\mathcal C_c}u}{\norm{\sum_{u\in\mathcal C_c}u}_2}.
 \label{eq:medoid}
\end{equation}
这里 $\argmax$ 返回成员向量，而非其索引。若平均方向范数过小，应保留可用代表或重新初始化。该成员代表与传统精确 K-means 均值更新不同，因此不直接套用经典 K-means 目标函数单调下降的证明。

Refine 阶段对近重复簇进行合并，对方向凝聚度不足的簇进行拆分。实验采用合并阈值0.9 和拆分阈值0.6：代表余弦相似度过高时合并，成员与代表平均相似度不足时使用二分聚类并继续局部调整。最大细化轮数限制离线开销。上述阈值是实验配置，不是保证任意数据达到目标召回的通用常数。

\begin{algorithm}[htbp]\caption{方向敏感聚类与簇结构细化}\label{alg:cluster}
\begin{algorithmic}[1]
\REQUIRE 稀疏向量集合，初始簇数，合并与拆分阈值
\ENSURE 簇集合、簇代表及成员索引
\STATE 归一化非零向量，按式\eqref{eq:spherical-init}初始化代表
\REPEAT
\STATE 按最大方向相似度分配成员
\STATE 按式\eqref{eq:medoid}更新实际成员代表
\UNTIL{达到收敛条件或迭代上限}
\STATE 合并代表高度相似的簇，并重新选择代表
\FOR{凝聚度不足的簇}
\STATE 进行二分聚类，重新分配成员并更新两个代表
\STATE 在细化预算内重复调整
\ENDFOR
\STATE 输出簇及其连续存储区间
\end{algorithmic}\end{algorithm}

\subsection{基于候选簇探测的选择性计算}
在线查询先计算与簇代表的相似度，再按探测预算选择候选簇。设选择集合为 $\mathcal P(x)$，则实际参与候选计算的行集合为 $\bigcup_{c\in\mathcal P(x)}\mathcal C_c$。由于各簇大小不同，簇数量比例不等于扫描向量比例；实验必须同时记录实际访问量，而不能仅以 $n_{\mathrm{probe}}/n_{\mathrm{list}}$ 估算运行时间。

硬件使用分布式簇探测，每个 ACPE 处理本地簇并接收 $sub\_nprobe$。原设计按 $\lceil n_{\mathrm{probe}}/32\rceil$ 分配每核预算，因此实际探测数还受取整和各核可用簇数限制。这种局部预算保证执行可扩展性，但与全局精确选择得分最高的 $n_{\mathrm{probe}}$ 个簇不同，需要通过端到端召回校准。

当探测预算增加，通常能够访问更多相关候选，但簇读取和内部计算成本也随之增加。对达到近全召回的设置，选择性收益可能显著缩小。因此，系统应面向目标召回进行参数选择，并保留接近全扫描时与基础加速器的比较。

\subsection{量化与非零元素剪枝优化}
聚类阶段使用浮点表示，以减少离线结构构建中的量化偏差；簇代表、候选与查询在检索阶段转换为 INT6。与第三章饱和裁剪形式不同，ParetoES 采用按最大绝对值缩放到 $[-31,31]$ 的对称量化：
\begin{equation}
 Q(v_j)=\round\left(\frac{31}{\max_t|v_t|}v_j\right).
\end{equation}
该范围选择避免额外的饱和裁剪，但舍入误差仍然存在。采用不同缩放粒度时还应保持比较分数的可比性，不能将独立缩放后的整数内积无条件视为原始内积排序。

本章沿用 ReSparse 的按幅值筛选思想，调整其阈值统计以适应高度稀疏数据。若均值包含大量零元素，全局平均绝对值可能很低，使多数行通过保护条件而无法有效剪枝。改用非零元素上的统计量可提高小值辨别能力，再结合聚类与探测预算在目标 Recall@100 下校准保留比例。该调整使第三章的元素级剪枝方法适配簇级选择性检索。

\section{面向选择性计算的 FPGA 加速架构}
\subsection{ACPE 整体架构}
ParetoES 在 Alveo U280 上配置32 个 ACPE，每个引擎绑定一个 HBM 伪通道，通过512 位接口读取本地数据。核心工作频率为225 MHz，32 路接口给出约460 GB/s 的理论传输能力；可持续有效带宽仍受具体访问模式和控制开销限制。
\thesisfig[0.67]{paretoes/Architecture.pdf}{ParetoES 的多 ACPE 与 HBM 通道绑定架构}{fig:pareto-arch}{%
各引擎在计算阶段独立处理本地簇，避免引擎间频繁同步。局部输出汇总后由主机形成最终候选结果。簇选择和簇内计算共用解码、乘法与归约流水，减少专用硬件重复。图\ref{fig:pareto-arch}所示架构通过相相同执行基础上的计算路径服务两个层级的筛选任务。
}


\subsection{候选簇筛选与计算流水线}
ACPE 启动时读取位于本地 HBM 数据区头部的簇代表，使用 Ultra-CSR 连续解码并计算其与查询的内积。每个传输包承载30 个 INT6 非零元素，查询向量通过15 份 URAM 副本提供30 个读取端口。归约后输出代表得分，交给局部排序单元确定待探测簇。
\thesisfig[1.00]{paretoes/ACPE.pdf}{ACPE 的簇筛选、地址映射与簇内计算流水线}{fig:pareto-acpe}{%
Mem Map 根据选中簇编号查找起始地址及范围，随后读取对应簇的数据块。簇间跳转具有随机性，簇内存储则按连续数据流组织。这种布局将原本全局散乱的访问限制在少量已选区域，且每次定位后可复用成熟的低比特稀疏流水。数据格式、索引结果与访存调度在簇粒度上保持一致，是选择性计算形成实际收益的关键。
}


\subsection{分布式 Top-K 排序与结果更新}
集中式 Bitonic 排序网络的比较器数量和深度随输入规模快速增长。对 $N=2^k$ 个输入，经典网络比较器数为 $Nk(k+1)/4$，深度为 $k(k+1)/2$。当 $N=512$ 时分别为11520 和45；当 $N=16$ 时为80 和10。32 个局部 Bitonic-16 总计2560 个比较器，显著小于单个512 输入网络。

分布式方案通过局部候选集合降低排序规模，其质量取决于局部容量对目标结果的覆盖程度。局部 Top-16 的并集至多包含512 项，与精确全局 Top-512 具有不同的选择语义。ACPE 进一步将 Top-16 更新拆为四个 Top-4 单元，通过轻量轮询减短关键路径，随后完成局部排序与编号--得分同步传递。

实验在相同候选集合上比较32 路局部 Top-16 与精确全局选择，对五个数据集的 $K\leq200$ 设置测得100\% Recall，$K=512$ 时约为0.899 至0.903。该实验隔离局部候选截断的影响，并不表示完整选择性检索在任意探测预算下均达到100\%召回。

\section{面向不规则计算的多核负载均衡方法}
\subsection{候选簇访问与负载分布特征}
选择性检索中，候选簇大小、热度和非零密度共同影响执行成本。仅按簇编号或数量均分，无法保证每个核的实际工作量接近。原论文在 Sp.Baidu 上用100 个查询统计簇级访问量，发现静态线性分配的 CV 为34.25\%，最大最小负载比为4.64，反映出明显执行拖尾\cite{paretoes2026}。

\thesisfig[1.00]{paretoes/baidu_hotspot.pdf}{簇热点分布与负载均衡效果}{fig:pareto-hotspot}{%
图\ref{fig:pareto-hotspot}说明热点分布与物理放置必须协同考虑。只将热点标识为“高频”并不足够，还需要避免多个热点聚集于同一通道；而对于不经常访问的大簇，也需保留对其实际工作量的统计。
}


\subsection{\texorpdfstring{\hbalance{}}{H2Balance} 负载均衡策略}
\hbalance{}采用热点画像、热点锚定和贪心分配三个步骤。先构建 $H[t,c]$，记录第 $t$ 个代表查询访问簇 $c$ 的向量工作量，并计算 $Q_c=\sum_tH[t,c]$。再将热度最高的至多32 个簇分别分配给不同核，防止主导负载的热点相互叠加。最后按累计负载将其余簇放置到当前最轻的核。
\begin{algorithm}[htbp]\caption{基于代表查询的分层热点均衡}\label{alg:hbalance}
\begin{algorithmic}[1]
\REQUIRE 簇访问统计 $H$，计算核数 $P$
\ENSURE 簇到核映射 $M$
\STATE $Q_c\leftarrow\sum_tH[t,c]$，将簇按 $Q_c$ 降序排列
\STATE 初始化各核负载 $L_p\leftarrow0$
\FOR{前 $\min(P,C)$ 个热点簇}
\STATE 分别锚定到不同核，更新对应 $L_p$
\ENDFOR
\FOR{其余簇 $c$}
\STATE $p^*\leftarrow\argmin_pL_p$
\STATE $M(c)\leftarrow p^*$，$L_{p^*}\leftarrow L_{p^*}+Q_c$
\ENDFOR
\STATE 按映射共同放置簇代表及成员数据
\end{algorithmic}\end{algorithm}

该策略依据历史查询统计离线生成映射，在线查询沿用已有的数据放置。若查询分布显著变化，需要重新画像并更新映射；更新本身涉及数据搬移，其频率应由收益和维护成本共同决定。按访问向量数量统计也可能忽略簇间非零密度差异，部署时可进一步验证统计量与真实传输字节数或执行时间的相关性。

\subsection{数据布局与执行调度协同优化}
完成映射后，将簇代表、地址表和成员数据放置在相应 ACPE 绑定的 HBM 通道。这样，逻辑上均衡的簇任务能够由本地数据流执行，避免跨通道读取重新引入瓶颈。运行时只需选择本地簇并顺序处理其数据区间，映射复杂度留在离线阶段。

负载均衡需要同时评估画像集合和未见查询。画像集合上接近完美的平均均衡，可能反映对已知分布的充分拟合；在未见查询或分布偏移下仍能降低 CV 与峰值负载，才能说明其具有一定鲁棒性。下一节分别考察画像查询、未见查询和受控分布偏移条件下的负载与性能。

\section{实验评估与分析}
\subsection{实验环境与评估方法}
本章依据 ParetoES 研究的统一实验配置进行评估\cite{paretoes2026}。CPU 为 Intel Xeon Gold 5117 平台，配置56 个 vCPU、2.00 GHz 及256 GB 内存；GPU 使用 A100 40 GB 及 V100-SXM2 32 GB，CUDA 12.1。Faiss 版本为1.7.2，采用内积 IndexIVFFlat 及 IndexFlatIP 量化器，另以\texttt{sparse\_dot\_topn}作为精确全量参考。FPGA 实现使用 Vitis HLS 2023.2 及 Alveo U280，对比 AccelES 与既有 FPGA32 方案。

\begin{table}[htbp]\centering
\caption{ParetoES 数据规模与离线预处理时间}\label{tab:pareto-data}
\begin{tabular}{lrrrrr}\toprule
数据集&行数&维度&密度&索引构建（s）&总预处理（s）\\\midrule
Sp.Baidu&640000&900&14.02\%&71&160\\
Sp.CC\_zh&2000000&900&2.12\%&92&165\\
Sp.GloVe&400000&500&13.83\%&37&86\\
Sp.Wiki News&1000000&900&3.04\%&76&144\\
Sp.10M&10000000&524&0.72\%&363&680\\\bottomrule
\end{tabular}
\end{table}

表\ref{tab:pareto-data}包括四个真实语料派生数据集及一个人工千万行数据集。Sp.10M 包含524 维特征，密度为0.72\%，其构造参数与第三章的人工10M 矩阵不同。索引构建在 A100 上完成，总预处理还包括量化、ReSparse 及均衡。在线 QPS 不包含一次性构建时间，因此适用于可摊销索引成本的多查询场景。

评估分为三类：同目标 Recall 下的跨平台 QPS 比较、与固定全量计算加速器的比较，以及固定架构下逐项启用算法和系统模块的消融。对 Recall@100 位于0.8 至1.0 区间的曲线汇总值与单点最大加速比分别取得，不将两种统计量相互替代。

\subsection{计算质量与吞吐率权衡}
图\ref{fig:pareto-performance}和图\ref{fig:pareto-performance10}分别展示五个数据集上的 Recall@100 与 Recall@10 曲线。探测预算增加时，系统访问更多候选，吞吐通常下降；不同聚类方法在相同扫描预算下的召回差异，决定了达到目标召回所需的实际工作量。完整 ParetoES 相对 A100 的加速范围为1.54 至79.34 倍，所报告几何平均为14.77 倍；相对 CPU 的范围为28.95 至540.73 倍。这些结果反映所评估数据集与召回区间内的端到端收益。
\thesisfig[0.86]{paretoes/performance.pdf}{Recall@100 目标下的端到端检索吞吐}{fig:pareto-performance}{%
以 Sp.Baidu 为例，表\ref{tab:pareto-qps}在接近的召回目标下列出 QPS。ParetoES 在约0.8 召回时达到4761.9 QPS，而接近1.0 召回时为1851.9 QPS。质量目标改变后吞吐并非保持不变，说明选择性计算将应用容忍度转化为实际扫描量的调节空间。
}

\thesisfig[.96]{paretoes/performance10.pdf}{Recall@10 目标下的端到端检索吞吐}{fig:pareto-performance10}{%
图\ref{fig:pareto-performance10}给出 Recall@10 下的性能比较，子图依次对应 Sp.Baidu、Sp.CC\_zh、Sp.GloVe、Sp.Wiki News 和 Sp.10M。与 Recall@100 相比，输出规模和排序边界发生变化，各系统达到相同召回目标所需的探测预算也相应改变。因此，两个指标分别评价，性能曲线均以各自的精确 Top-K 集合作为参考。
}

\begin{table}[htbp]\centering
\caption{Sp.Baidu 在匹配召回水平下的 QPS}\label{tab:pareto-qps}
\begin{tabular}{lrrr}\toprule
系统&Recall 约0.8&Recall 约0.9&Recall 约1.0\\\midrule
ParetoES&4761.9&2857.1&1851.9\\
Faiss-A100&1172.5&834.9&558.9\\
Faiss-V100&806.5&567.9&357.7\\
Faiss-CPU&45.6&30.3&22.9\\
AccelES&1818.2&1818.2&1818.2\\
FPGA32&425.5&425.5&425.5\\\bottomrule
\end{tabular}
\end{table}

对应 ParetoES 探测配置为128/398、224/398 和384/398，Faiss 配置为180/400、240/400 和400/400。相同召回并不要求相同探测数，聚类质量本身属于系统收益来源。若需隔离架构效果，应进一步固定聚类方法，不能将该表所有差异归因于 FPGA 硬件。

\subsection{与全量计算加速器的性能比较}
在 Recall@100 为0.8 至1.0 的评估区间内，ParetoES 相对 AccelES 和 FPGA32 的平均吞吐提升分别为2.27 和7.61 倍；接近全召回时，相应平均提升缩小为1.21 和3.89 倍。Sp.Baidu 上接近全召回的 ParetoES 仅比 AccelES 高约1.02 倍，说明范围裁剪优势接近消失后，簇定位开销会抵消部分收益。

上述结果说明，ParetoES 通过可配置的探测预算提供质量--吞吐率权衡，其收益随召回目标和数据结构变化。允许适度召回损失时，候选范围缩减带来更显著的优势。对需要极高候选覆盖且数据库结构不易分离的任务，AccelES 仍具有流程较简单、行为较稳定的价值。

\subsection{关键技术消融与负载均衡分析}
增强 ReSparse 在五个数据集上报告平均37.41\%和最高61.25\%的非零剪枝比例。原始 ReSparse 在 Sp.10M 上只删除2.34\%，改进后为11.48\%，说明非零统计阈值能够改善极高稀疏度下的剪枝能力。图\ref{fig:pareto-nnz}展示相应非零数变化，质量仍需结合簇选择后的最终召回判断。
\thesisfig[0.85]{paretoes/nnz_absolute_comparison.pdf}{量化与增强 ReSparse 对非零元素数量的影响}{fig:pareto-nnz}{%
图\ref{fig:pareto-nnz}表明，非零元素的减少幅度会随数据集和剪枝策略变化。因而，元素数量是解释访问成本的中间指标，仍需与最终召回率和实际扫描范围结合判断其有效性。
}


\thesisfig[0.62]{paretoes/roofline_ablationStack.pdf}{ParetoES 优化组合的吞吐变化}{fig:pareto-ablation}{%
Sp.CC\_zh 上的优化栈从 K-means 配置的3307 QPS，依次提高到聚类细化后的3622 QPS、加入\hbalance{}后的4065 QPS，以及完整配置的5869 QPS。图\ref{fig:pareto-ablation}展示组合收益。该序列显示各模块组合后获得持续改善，相邻配置的增幅体现模块在已有优化基础上的增量贡献，包含模块之间的交互作用。
}


\hbalance{}在画像查询集合上将 Sp.Baidu 负载 CV 由34.25\%降至约0.01\%；在100 个未见查询上降至5.16\%，最大最小比降至1.24，并在 Recall@100 为0.8 时带来37.52\%吞吐提升。对查询加入0.2 正偏置后再归一化的受控偏移实验中，CV 由36.27\%降至12.25\%，最大最小比由4.48 降至1.63。这些结果说明历史映射在一定分布变化下仍有帮助，但分布偏移会削弱均衡效果，表明映射有效性依赖画像与实际查询分布的接近程度。

实验还固定 K-means 基线以隔离架构差异，在匹配 Recall 区间下报告相对 A100 和 CPU 分别7.9 和194.5 倍 QPS。将 ReSparse 应用于 Faiss 后，CPU 平均提高约14\%，GPU 却下降约7\%。这说明元素减少是否形成加速，取决于稀疏执行结构能否有效利用剩余工作量。

\subsection{能耗与硬件资源开销分析}
ParetoES 利用 INT6 与 DSP 复用降低算术资源成本，同时为30 路查询读取保留较多 URAM 副本。分布式微型排序增加局部逻辑，但避免集中式大规模排序网络。通过将 Top-16 更新分解为多个小单元，核心达到225 MHz。测得功耗为59 W，在相应资源评估配置下能效为1.38 GOPS/W。

吞吐提升能否转化为单查询能耗下降，可用 $E_{\mathrm{query}}=P\,T_{\mathrm{query}}$ 分析。即使即时功耗略有增加，缩短处理时间仍可能降低单次能耗。但应保持功耗测量范围一致，并区分板卡与整机功耗。本章侧重运行功耗与单查询能耗，部署成本还取决于设备价格、利用率和服务周期。
\thesisfig[0.67]{paretoes/P99_latency.pdf}{Sp.GloVe 不同召回目标下的 P99 查询延迟}{fig:pareto-latency}{%
图\ref{fig:pareto-latency}关注尾延迟，而不仅是平均 QPS。在 Sp.GloVe 高召回区间，实验测得相对 Faiss-CPU 与 Faiss-GPU 的平均 P99 加速分别为179.6 和5.5 倍。该结果与簇选择减少工作量、热点均衡减轻执行拖尾相一致；若要区分二者对尾延迟的独立贡献，还需要在相同查询序列上分别启用模块进行评估。
}


% !TeX root = ../../main.tex
\section{选择性执行的成本结构与稳健性讨论}
\subsection{簇代表成本与成员扫描成本}
选择性检索增加了簇代表比较。设簇数为 $C$，代表矩阵非零数为 $z_c$，选中成员的非零数为 $z_s$，则在线稀疏工作量至少包含 $z_c+z_s$。如果簇代表趋于稠密或簇数过多，即使成员扫描量下降，代表计算也可能成为不可忽略的固定成本。

采用实际稀疏成员作为代表能够约束这一成本，但可能牺牲均值代表的拟合能力。选择代表时需联合观察代表非零数、簇内方向凝聚度及达到目标召回的探测数。仅比较聚类目标函数不足以确定检索系统性能，因为代表稀疏性与在线执行成本并不直接出现在普通 K-means 损失中。

\begin{table}[htbp]\centering\caption{候选簇粒度对系统成本的双向影响}
\begin{tabularx}{\textwidth}{p{2.6cm}XX}\toprule
变化&潜在收益&潜在代价\\\midrule
簇数量增加&更精细的范围选择&代表比较和地址表增长\\
簇数量减少&减少索引与探测固定开销&单簇扫描范围变大\\
代表更稀疏&降低代表访存&可能降低方向代表性\\
簇内更紧凑&相关候选更集中&离线细化成本增加\\
簇大小更均匀&降低核间极端负载&可能偏离相似度最优划分\\\bottomrule
\end{tabularx}\end{table}

\subsection{簇级召回与元素剪枝的交互}
簇级选择排除整组候选，元素剪枝改变保留候选的近似得分。两种误差作用位置不同：被排除簇中的真实 Top-K 无法通过提高簇内精度恢复，而被量化扰乱的候选可能通过更多预算或原始得分重算改善。因此，调参应先确认簇范围有足够覆盖，再讨论簇内表示精度。

设所探测簇包含的真实 Top-K 比例为 $R_c$，在这些真实候选中经过局部筛选仍保留的条件比例为 $R_l$，则完整候选保持可写为 $R_cR_l$。这里是基于同一查询集合事件划分的集合计数关系，不是假设两种误差独立。对跨查询平均值，平均乘积通常不等于两个平均值的乘积，故不应通过分别测得的平均召回直接相乘估算端到端结果。

这个分解可以指导消融：一组实验以原始得分在选中簇内精确排序，测量范围选择的上限；另一组固定同一候选范围，比较 INT6 和 ReSparse 后的结果。最后再启用分布式候选维护，检查硬件近似是否增加损失。完整结果仍以端到端测量为准。

\subsection{探测预算取整与核内选择}
32 核系统将全局预算分解到本地执行。若 $n_{\mathrm{probe}}$不是32 的倍数，统一向上取整会使实际预算最多多出31 个簇；若某核簇数不足，还需截断其预算。因此，命令参数并不必然等于真正访问簇数。实验应记录实际探测簇、实际候选数与实际非零数三个量。

\begin{table}[htbp]\centering\caption{探测预算的三个层次}
\begin{tabularx}{\textwidth}{p{2.5cm}XX}\toprule
指标&含义&为何不能互相替代\\\midrule
请求探测数&软件传入的预算&受到分布式分解和取整影响\\
实际簇数&所有核真正访问簇的数量&不同簇包含的成员不同\\
实际行数&参与相似度计算的向量数量&不同向量非零数不同\\
实际非零数&需要传输和乘加的记录数量&还未包含索引和事务开销\\\bottomrule
\end{tabularx}\end{table}

局部预算也可能造成相关簇集中区域被限制。增加每核预算能够缓解，但会在不相关核上增加访问。更灵活的预算分配可以基于代表得分或历史需求调整，然而需要额外协调信息，与完全独立 ACPE 执行存在取舍。本文采用固定分解获得较简单控制路径，并通过实测召回评价其结果。

\subsection{热点画像的时间尺度}
热点画像是对查询分布的统计估计。画像查询过少，容易受偶然候选影响；过长的历史窗口则可能掩盖近期变化。对于静态数据但查询意图变化的服务，簇结构仍有效而热点映射可能失效，因此索引重建与映射更新应视为不同层次的维护操作。

设新映射预计每次查询节省 $\Delta t$，迁移成本为 $T_m$，预测在下一次明显分布变化前还能处理 $N$个查询，则需要 $N\Delta t>T_m$才能获得总时间收益。这里的节省还应包含迁移期间的干扰，而不是只比较迁移前后的稳定状态 QPS。

原论文的受控正偏置实验验证一种特定分布变化，不代表所有真实漂移。后续可区分查询主题变化、范数变化、稀疏度变化与突发热点，分别检查 CV、最大核时间和尾延迟。只有把漂移类型与质量保持一起记录，才能判断映射是否真正稳健。

\subsection{长尾簇与不可分割任务}
贪心分配能均衡多个簇的累计工作，但若一个簇本身的工作量远超平均值，即使独占一个核也会形成瓶颈。将任务视为不可分割单元时，理论完成时间至少为最大单簇负载与平均负载两者的最大值。这一界解释了仅靠重新放置无法解决所有热点。

拆分热点簇能够增加调度自由度，却可能改变代表集合和探测语义；复制热点数据则需要同步更新并消耗容量。两者都需要重新定义结果合并和去重规则。如果复制后同一候选可由多个核返回，应在最终 Top-K 前去重，防止重复 ID 挤占结果容量。

\subsection{性能曲线与前沿比较方法}
对每个系统，通过多个探测预算获得 $(R,Q)$点，其中 $R$为召回、$Q$为 QPS。在相同或更高召回下仍具有更高吞吐的点，支配被比较点。可移除被支配点构成观测前沿，但这只覆盖实际测量的参数空间，不等价于全部配置的最优边界。

如果目标召回位于两个测量点之间，可以报告达到或超过目标的最小预算点，或明确说明使用插值。不同系统若使用不同插值方法，可能夸大差异。对固定全量系统，其质量和吞吐是固定点或近似水平参考，不应将其误画成能够连续调节召回的算法。

还应区分算术平均 QPS、几何平均加速比和最大加速比。前者受绝对吞吐规模影响，后者衡量相对比值，最大值只描述最有利案例。本文分别报告上述统计量，并明确其对应的数据集合与质量区间。

\subsection{尾延迟、能耗与在线服务}
平均 QPS 无法直接描述尾延迟。同步多核查询由最晚完成的核决定，偶发热点会使少量查询耗时较长；系统排队又可能进一步放大尾部。在离线单查询测量中观察到的 P99 改善，不自动等同于固定到达率服务下的 P99 改善。

在线服务评估需要给出到达过程、并发数和预热方式，并区分设备执行时间与排队时间。能耗应对整个测量窗口积分，而不是仅使用标称功耗乘以一个最短内核时间。若 CPU 负责候选合并或索引维护，其成本也应纳入相应的系统边界。

\subsection{高维数据的适配边界}
当前 ACPE 继承10 位列索引和多份片上查询缓存，适合论文中的维度范围。对数万维稀疏词项向量，列索引变宽会减少每包非零数量，完整复制查询也可能不可行。可考虑分块查询缓存或稀疏查询索引，但会引入额外间接访问和缓存未命中。

因此，将本系统拓展到高维词项检索需重新评估数据格式、簇代表非零度和查询缓存路径，而不是只提高一个维度参数。该边界也说明实验材料中的经典词向量稀疏化数据与现代神经稀疏检索数据应分别取得，保持论证范围清楚。

% !TeX root = ../../main.tex
\section{结构选择与多核执行的分析性示例}
\subsection{方向相似性与幅值的不同作用}
考虑三个二维向量 $a=(1,0)$、$b=(10,0)$、$c=(1,1)$。欧氏距离满足 $\norm{a-b}=9$、$\norm{a-c}=1$，但归一化后的余弦相似度分别为1 和 $1/\sqrt2$。若希望聚类优先反映方向，单纯依据原始欧氏距离可能把同方向但不同范数的向量分开。球面聚类通过归一化强调方向关系，与这一需求相适应。

该例用于说明度量差异，不意味着数据库行范数可以在最终内积中被任意删除。内积满足 $a^{\mathsf T}q=\norm a\norm q\cos\theta$，幅值仍然影响最终排名。聚类阶段的规范化与最终评分的数据表示必须区分：前者帮助形成可探测区域，后者仍需遵守原得分定义。否则，聚类优化就会无意中改变检索问题本身。

稀疏数据还具有支撑集合差异。两个向量可能只有少量坐标重合，均值代表则可能合并多成员支撑而变得更稠密。使用真实成员作为代表有助于维持稀疏性和统一计算路径，但代表并不等于所有成员的得分上界。第4 章的代表筛选因此是一种近似范围决策，需要由召回实验验证，而不能视为可证明安全的界限剪枝。

\subsection{簇内一致性与代表可辨识性}
好的聚类不仅应使成员彼此接近，还应使查询能够通过代表分数识别相关簇。两种目标存在联系，但不完全等价。假设某簇含大量常见方向和少量边缘方向，代表可能更好地反映多数成员；当查询与边缘方向相关时，簇整体代表分数仍可能较低，从而遗漏少数高分成员。

因此，评价聚类不能只报告目标函数下降。还需要观察目标 Top-K 分布于多少簇、相关簇在代表排序中的位置，以及探测预算增加时累计召回的变化。若真实 Top-K 集中在少量容易识别的簇，选择性计算具有明显机会；若目标分散且代表不敏感，则即便簇内距离较小，也可能需要较高探测比例。

可将一次查询真实 Top-K 所在的簇集合记为 $\mathcal G_K(q)$，被探测簇集合记为 $\mathcal P(q)$。只要存在某个目标行所属簇不在 $\mathcal P(q)$ 中，该行就不可能由后续簇内重算恢复。对这一层而言，增加候选精排精度无济于事，需要扩大探测范围或改善代表及聚类结构。

\subsection{簇选择和簇内候选的联合影响}
设真实目标集合为 $\mathcal T$，被探测区域包含的行集合为 $\mathcal V$，精排前候选为 $\mathcal C\subseteq\mathcal V$。可定义区域覆盖率与条件保留率：
\begin{equation}
 r_{\mathrm{region}}=\frac{|\mathcal T\cap\mathcal V|}{|\mathcal T|},\qquad
 r_{\mathrm{keep}}=\frac{|\mathcal T\cap\mathcal C|}{|\mathcal T\cap\mathcal V|}.
\end{equation}
在分母非零时，二者乘积等于候选覆盖率。这是集合计数的恒等关系，不要求误差事件独立；当分母为零时，最终覆盖为零，条件比率则应记为未定义，不能强行代入。

例如真实 Top-10 中有8 行进入探测区域，簇内低比特筛选保留其中7 行，则区域覆盖为0.8，条件保留为 $7/8$，最终候选覆盖为0.7。这个人为构造的例子说明，只提高簇内保留率最多把覆盖恢复到0.8；要进一步提高，需要改变探测区域。它同样不意味着对多条查询分别求平均后仍可直接相乘，平均方式必须保留权重或逐查询统计。

\begin{table}[htbp]\centering\caption{两层选择的诊断结果与调整方向}
\begin{tabularx}{\textwidth}{p{3.4cm}XX}\toprule
观察结果&主要受限环节&优先考察的参数\\\midrule
区域覆盖低、条件保留高&相关簇未被探测&代表质量、聚类与探测预算\\
区域覆盖高、条件保留低&簇内近似或局部截断&剪枝、量化、局部候选容量\\
两者均高但吞吐低&执行效率或固定开销&代表扫描、地址跳转、均衡\\
两者均低&多个机制同时限制&分层对照与联合参数搜索\\\bottomrule
\end{tabularx}\end{table}

\subsection{从热点统计到实际完成时间}
热点均衡需要估计每个簇被访问的频率以及每次访问的成本。若簇 $j$ 的访问概率为 $p_j$、有效工作量为 $w_j$，则 $p_jw_j$可作为平均贡献的近似指标。只按簇大小分配会忽略访问频率，只按频率分配又可能忽略大小差异。实际成本还受到格式、访存和排序影响，因此该乘积是映射启发信息，不是精确服务时间。

考虑四个簇，其每次访问工作量为8、8、2、2 个任意单位，访问概率为0.8、0.2、0.8、0.2。期望工作量分别为6.4、1.6、1.6、0.4。若两个核分别放置前两个与后两个簇，期望负载为8 和2；若将第1、4 簇放在一核，第2、3 簇放在另一核，则变为6.8 和3.2。该例说明热点参与映射可以改善平均不均衡，但最大热点本身仍限制可达到的均衡程度。

这些数字是用于解释的统计量，并非 H²Balance 的实测结果。更重要的是，期望负载均衡并不保证每条查询均衡。如果若干热簇总是被同一查询共同访问，把它们放在同一核可能产生突发执行拖尾，即便长期平均相近。完整分析需要查询级访问关联，而不能只看各簇的边际频率。

本研究利用代表性查询估计负载，并通过未见查询与受控分布偏移检验映射的适应性。后续若引入更细粒度关联模型，应把统计成本、重新布局成本和映射计算成本纳入评价，避免为了改善少量尾部请求付出过高的更新代价。

\subsection{探测预算取整与小预算区域}
多核系统通常要求每核执行整数个探测任务。若全局预算为 $n_{\mathrm{probe}}$，以 $P$ 个核均分时采用向上取整，则实际可分配槽位为
\begin{equation}
 P\left\lceil n_{\mathrm{probe}}/P\right\rceil.
\end{equation}
当预算远大于核数时，取整误差相对较小；当预算与核数同量级时，参数增加一点就可能跨越一个整轮并行任务，造成工作量阶跃。因此，探测预算与延迟未必呈平滑线性关系。

还需区分槽位数和实际访问簇数。如果某核可用簇不足，或目录存在空簇，实际读取可能小于槽位上限；如果实现要求每核至少探测一个簇，小预算则受到并行组织下限约束。复现实验应记录最终实际扫描行数、非零数和读取字节数，用它们解释性能，不能只使用名义探测参数。

\subsection{质量约束下的性能工作点}
选择性检索应以质量与吞吐成对报告。设配置 $a$具有质量 $r_a$、吞吐 $q_a$，若另一配置 $b$同时满足 $r_b\geq r_a$、$q_b\geq q_a$，且至少一项严格更好，则 $a$被 $b$支配。非支配配置构成经验上的 Pareto 前沿。这里的前沿只覆盖已经测量的配置集合，不保证已经找到所有可能算法或参数中的全局最优。

在目标召回 $r_{\min}$下，可从已测配置中选择满足 $r\geq r_{\min}$且吞吐最高的点。若两个系统只在不同召回位置有数据，应增加匹配质量的测量，或明确说明插值规则。直接把低召回点与高召回点的吞吐作比，会把质量放宽误算为纯执行优势。

\begin{table}[htbp]\centering\caption{选择性检索工作点的报告内容}
\begin{tabularx}{\textwidth}{p{2.8cm}X}\toprule
项目&指标含义\\\midrule
质量目标&Recall@K 中的 K、精确参考集合、并列规则\\
范围参数&总簇数、名义探测数、实际探测与扫描比例\\
数值参数&位宽、非零保留策略、候选容量\\
性能测量&吞吐、单查询延迟、尾延迟及批量方式\\
资源代价&板卡资源、索引容量、构建与重布局时间\\
统计方式&查询集合、重复次数、聚合与异常处理\\\bottomrule
\end{tabularx}\end{table}

接近全召回时，系统需要访问更多簇，选择性带来的工作量优势减小，而代表扫描和簇目录处理仍然存在。实验中相对 AccelES 的收益随质量目标变化，正符合这一机制。不能把某一宽松工作点的高加速比推广到全部召回范围，也不能据此否定高召回下仍可能存在的布局与排序收益。

\subsection{静态布局与在线适应之间的取舍}
静态热点布局将统计和迁移开销放在离线，在线控制简单，适合分布相对稳定的查询。在线重平衡能响应分布变化，但需要监控、决策与数据移动；迁移期间的双份目录和查询一致性也会增加工程复杂度。因此，是否更新应取决于预期节省能否覆盖迁移成本，而非检测到任何偏移就立即重建。

对更新较频繁的数据，可考虑把静态主体与小型增量区域分开搜索，再合并候选。这是一种可能的后续实现方向，本文尚未测量其资源和质量效果。它要求重新分析两个区域的候选容量以及总查询成本，不能简单认为增量区域足够小就没有影响。

本章的结构选择、硬件探测和热点均衡共同作用于实际工作量与关键路径。理解它们之间的联系，有助于解释为什么仅增加聚类质量或仅减少扫描比例不一定得到相同比例加速，也有助于明确系统在高维、动态更新和分布偏移条件下需要进一步验证的部分。

\section{本章小结}
本章针对全量低比特扫描仍存在的计算范围冗余，提出方向敏感聚类、候选簇探测和相应 FPGA 选择性执行机制。ACPE、分布式微型排序与\hbalance{}共同解决簇级选择引入的地址跳转、比较资源和热点负载问题。在继承第三章低位宽数据流的基础上，该方法将优化重点推进到实际访问范围，并以召回率--吞吐曲线评价收益。其质量要求仍属于检索候选层面，下一章将研究拓展到误差影响后续状态的迭代稀疏计算。
